\section{The End That Is Offered}

What revelation asserts, at the precise joint where the ascent went silent, is that the ceiling has been opened from above. The First Vatican Council---the same constitution this series has quoted on free creation and on total dependence---states it in a single Latin clause: revelation is necessary because ``God, out of His infinite goodness, ordered man to a supernatural end'' (\latin{Deus ex infinita bonitate sua ordinavit hominem ad finem supernaturalem}), namely ``to share in divine goods which altogether surpass the understanding of the human mind''; and the council seals the clause with the sentence Scripture keeps at that door: ``eye hath not seen, nor ear heard, neither hath it entered into the heart of man, what things God hath prepared for them that love him.''\endnote{First Vatican Council, dogmatic constitution \emph{Dei Filius} (24 April 1870), ch.~2 (\latin{Deus ex infinita bonitate sua ordinavit hominem ad finem supernaturalem, ad participanda scilicet bona divina, quae humanae mentis intelligentiam omnino superant; siquidem oculus non vidit, nec auris audivit, nec in cor hominis ascendit, quae praeparavit Deus iis, qui diligunt illum}, citing 1 Cor 2:9). Latin checked 2026-07-25 at the Holy See web text; the English renderings here are a labeled working gloss governed by the quoted Latin, except that the scriptural clause is given in the Douay--Rheims wording of 1 Corinthians 2:9 (``That eye hath not seen, nor ear heard, neither hath it entered into the heart of man, what things God hath prepared for them that love him''), checked 2026-07-25 at drbo.org. Note the council's own restraint: the supernatural end flows from God's goodness (\latin{ex infinita bonitate sua}), not from any exigency of nature. Part I quoted ch.~1 of this constitution (creation \latin{ex nihilo}, ``not for increasing his own beatitude''); Part II quoted ch.~3 (total dependence grounding the obedience of faith).} \emph{Ordered to}: the grammar of ends again, now on the giver's side of the relation. The creature was not merely built with a reach exceeding its arms; it was aimed, deliberately, at the very thing its arms cannot reach---aimed by the one being who can close the gap, in an act the council is careful to root in goodness, not in debt. Nothing in the creature's nature extorted this. The gift beneath the creature's being, which Part I uncovered, turns out to have a second story: the same liberality that gives existence has appointed, for the existing creature, an end above every created order.

The Scriptures name the end with a bluntness the metaphysics can only envy. Paul, to the Corinthians: ``We see now through a glass in a dark manner; but then face to face. Now I know in part; but then I shall know even as I am known.''\endnote{1 Corinthians 13:12, Douay--Rheims (Challoner revision); Scripture quotations in this essay are from the public-domain Douay--Rheims translation, wording checked 2026-07-25 at the drbo.org text. ``Through a glass in a dark manner'' renders \latin{per speculum in aenigmate}: the mirror-and-riddle knowledge of God from creatures is exactly the mode of knowing this series' first two parts exercised; the verse's own contrast term is the face.} John, in his first letter, in two sentences that hold the whole structure of this series: ``Behold what manner of charity the Father hath bestowed upon us, that we should be called, and should be the sons of God\ldots{} Dearly beloved, we are now the sons of God; and it hath not yet appeared what we shall be. We know, that, when he shall appear, we shall be like to him: because we shall see him as he is.''\endnote{1 John 3:1--2, Douay--Rheims, checked 2026-07-25 as above. The exegetical use made here is modest and textual: the verses assert a present filiation (``we are \emph{now} the sons of God''), a real but undisclosed future (``it hath not yet appeared''), and a likeness effected by vision (``we shall be like to him: because we shall see him as he is''). Aquinas quotes verse 2 at exactly the joint where this essay uses it, ST I, q.~12, a.~5, cited below at the light of glory.} And behind both, the old psalm this series' second part learned to read---the psalmist who asked what he should render---has a companion verse that now sounds like the whole human ascent compressed to a cry: ``My heart hath said to thee: My face hath sought thee: thy face, O Lord, will I still seek.''\endnote{Psalm 26:8, Douay--Rheims (Vulgate numbering), checked 2026-07-25 as above. The verse is cited as the psalter's own register of the desire analyzed in the preceding sections---the heart seeking not a doctrine but a face---and not as an exegetical claim about the psalm's original setting.} A face sought; a face promised; a seeing that transforms the seer into the likeness of the seen. Revelation does not correct the ascent's finding that the end is vision. It confirms it past all daring, and then adds the one thing the ascent could not: \emph{it is offered}.

How such seeing is possible for a creature, the tradition explains with a doctrine that repeats, at the summit, the exact structure Part I found at the foundation. The created intellect sees the essence of God only when ``God by His grace unites Himself to the created intellect, as an object made intelligible to it''; and for the creature to bear that union, something must be given to it---a strengthening the tradition calls the \emph{light of glory} (\latin{lumen gloriae}), by which, Aquinas writes, ``the blessed are made `deiform'---i.e.\ like to God, according to the saying: `When He shall appear we shall be like to Him, and we shall see Him as He is.'\,''\endnote{\emph{Summa theologiae} I, q.~12, a.~4, corpus (the grace-union), with q.~12, a.~5, corpus (the created light by which ``the blessed are made \emph{deiform}''---\latin{efficiuntur deiformes, idest Deo similes}---quoting 1 John 3:2); the phrase \latin{lumen gloriae}, the light of glory strengthening the intellect to see God (\latin{confortans intellectum ad videndum Deum}), stands in q.~12, a.~2, corpus. All checked 2026-07-25, English at New Advent, Latin at Corpus Thomisticum. The light of glory is a created disposition of the seer, not a created image of the seen: what is seen is the divine essence itself, not a likeness---the point the dogmatic definition below makes with the phrase ``no creature mediating.''} Attend to the shape of that sentence. The creature does not see by finally flexing hard enough; it sees because even its seeing is received---eyes lifted by the very light they look at, as the air of Part I was bright with a brightness not its own. At the top of everything, the lit-air structure does not give way to some achievement at last the creature's own. It intensifies: the vision of God is the gift of God, given \emph{in} the seer, and the creature's last act is like its first existence---a reception all the way down. What the vertigo disclosed at the bottom of being is the very mechanism of glory.

The Church has fixed this hope in her most deliberate register, and the fixing has a date and an occasion. In 1336, ending a controversy his own predecessor had stirred, Benedict XII defined that the purified souls of the saints, even now, before the resurrection and the general judgment, ``have seen and do see the divine essence with an intuitive vision, and even face to face, without the mediation of any creature''---\latin{visione intuitiva et etiam faciali, nulla mediante creatura}---``the divine essence immediately showing itself to them, plainly, clearly and openly'' (\latin{divina essentia immediate se nude, clare et aperte eis ostendente}); and that so seeing they possess it in joy, ``truly blessed,'' having ``eternal life and rest.''\endnote{Benedict XII, apostolic constitution \emph{Benedictus Deus} (29 January 1336), Latin text as received in Denzinger--Hünermann, \emph{Enchiridion symbolorum}, nn.~1000--1002 (\latin{vident divinam essentiam visione intuitiva et etiam faciali, nulla mediante creatura in ratione obiecti visi se habente, sed divina essentia immediate se nude, clare et aperte eis ostendente\ldots{} sunt vere beatae et habent vitam et requiem aeternam}); Latin phrases checked 2026-07-25 at a public bilingual transcription of DH 1000--1002, and the English of the central definition quoted as given in \emph{Catechism of the Catholic Church} 1023, official Vatican English web text, checked 2026-07-25 (``have seen and do see the divine essence with an intuitive vision, and even face to face, without the mediation of any creature''). The remaining English renderings are a labeled working gloss governed by the quoted Latin. The definition concerns the souls of the purified dead prior to the resurrection; questions beyond its own scope lie outside this essay.} A dogmatic definition is the Church speaking with maximum deliberateness, and it is worth hearing what this one is \emph{for}. Not speculation satisfied; a promise guaranteed. \emph{No creature mediating}: the seeing will not be one more inference from effects, one more song argued back to its singer. \emph{Nude, clare et aperte}: the God whom Part I could approach only through the analogy of lamps and springs, and whom the apophatic masters warned every concept would misdescribe, will show \emph{Himself}---and the initiative, note the participle, is His: \latin{ostendente}, showing. The one-sided relation that made the vertigo dizzying---the source that receives nothing, on whom nothing can be inflicted---was never, it turns out, a closed door. It was the silence of someone intending to be seen rather than deduced.

And with that, the arrest of Part I completes its long transformation. The creature hangs, now as ever, wholly suspended from what it cannot repay or affect. But the suspension has been given a direction and a terminus: the hands beneath the falling creature are carrying it \emph{toward the face}. What remains is to ask what such a promise makes of the relation in the meantime---what it is called when the one who cannot be repaid nevertheless communicates his own happiness to the debtor. The tradition's answer is the most startling single word in this series, and it comes from the mouth of the promise himself, at supper, on the night he was betrayed.
